\section{Compiling a changing observation parameter into PR syntax}
\label{sec:demand}

The finite-demand packet already gives the mathematical construction in
this section~\cite{repo-demand}. Its explicit outstanding task was to
construct the \code{PR} syntax and connect it to object arithmetic.
We reproduce the argument so that this report is self-contained, then
describe the new checked implementation. The distinction between the
inherited theorem and the new formalization is intentional.

\subsection{The equations that need compilation}

Let \(D:\N^2\to\N\), \(b:\N\to\N\), and
\(F:\N^3\to\N\) be fixed PR functions. We want
\begin{align}
 X(0,k)&=b(k),\label{eq:demand-zero}\\
 X(n+1,k)&=F\bigl(n,k,X(n,D(n,k))\bigr).
 \label{eq:demand-successor}
\end{align}
The recursive call changes \(k\). It is therefore not the ordinary PR
constructor with \(k\) held fixed. Treating \(X(n,\cdot)\) as an
accumulator of function type would also leave the first-order PR algebra.
Nevertheless, the future requests are determined by \(D\) before the
recursive answers are known. This gives a direct first-order solution.

For terminal data \((N,K)\), define
\begin{align}
 B(0;N,K)&=K,\nonumber\\
 B(j+1;N,K)&=D\bigl(N\dotminus(j+1),B(j;N,K)\bigr),
 \label{eq:backward}\\
 H(t;N,K)&=B(N\dotminus t;N,K).
 \label{eq:horizon}
\end{align}
Truncated subtraction makes the definitions total at every input.
Only \(t\le N\) is used in the correctness argument.
Then perform a forward recursion with the terminal data fixed:
\begin{align}
 V(0;N,K)&=b(H(0;N,K)),\nonumber\\
 V(t+1;N,K)&=F\bigl(t,H(t+1;N,K),V(t;N,K)\bigr),
 \label{eq:forward}\\
 X(N,K)&=V(N;N,K).
 \label{eq:compiled}
\end{align}
Both passes use ordinary primitive recursion. The forward pass may
recompute a demanded horizon instead of storing a whole table; either
choice has fixed PR syntax.

\begin{lemma}[Horizon identities]
\label{lem:horizon}
The definitions satisfy
\[
 H(N;N,K)=K,\qquad
 H(t;N,K)=D(t,H(t+1;N,K))\quad(t<N).
\]
They also satisfy the coherence identities
\begin{align}
 B(j+1;N+1,K)&=B(j;N,D(N,K)),\label{eq:back-coherent}\\
 H(t;N+1,K)&=H(t;N,D(N,K))\quad(t\le N).
 \label{eq:horizon-coherent}
\end{align}
These equations have quantifier-free induction proofs in \(\ipra\).
\end{lemma}

\begin{proof}
The terminal equation follows from \(N\dotminus N=0\) and the base
equation for \(B\). If \(t<N\), then
\[
 N\dotminus t=(N\dotminus(t+1))+1,
 \quad
 N\dotminus((N\dotminus(t+1))+1)=t.
\]
Substitute these into the successor equation for \(B\).

For \eqref{eq:back-coherent}, induct on \(j\). At \(j=0\), the left
side is \(D(N,K)\), which is the right base value. In the successor
case use the previous equality and the elementary identity
\[
 (N+1)\dotminus(j+2)=N\dotminus(j+1).
\]
This identity is valid even beyond the requested horizon. The resulting
equations have identical outer \(D\) arguments. Finally, for \(t\le N\),
\((N+1)\dotminus t=(N\dotminus t)+1\); substituting this value for
\(j+1\) in \eqref{eq:back-coherent} proves
\eqref{eq:horizon-coherent}. All numerical side conditions are Boolean
comparisons, and the induction formulas are equations or guarded equations.
\end{proof}

\begin{theorem}[Predetermined demand recursion]
\label{thm:demand}
The fixed PR function \(X\) of \eqref{eq:compiled} satisfies
\eqref{eq:demand-zero} and \eqref{eq:demand-successor}. Its construction
and the two equations have \(\ipra\) proofs for each fixed triple of
input PR definitions.
\end{theorem}

\begin{proof}
At \(N=0\), the horizon is \(K\) and the forward base is \(b(K)\).
For the successor equation compare the forward executions for
\((N+1,K)\) and \((N,D(N,K))\). Their initial values are equal by
\eqref{eq:horizon-coherent}. At every step \(t<N\), their horizon
values at \(t+1\) are equal by the same equation. Quantifier-free
induction on the guarded numerical equality
\[
 t\le N\ \longrightarrow\
 V(t;N+1,K)=V(t;N,D(N,K))
\]
therefore makes their values equal through time \(N\). The final step
of the longer execution is
\begin{align*}
 X(N+1,K)
 &=F(N,H(N+1;N+1,K),V(N;N+1,K))\\
 &=F(N,K,V(N;N,D(N,K)))\\
 &=F(N,K,X(N,D(N,K))).
\end{align*}
The middle equality is the terminal horizon equation and the guarded
forward equality just proved. No induction hypothesis quantifies over
all demands or over an infinite family of completed objects.
\end{proof}

\subsection{Correctness at a requested observation}

Suppose \(C(n,k)\) is bounded and \(\ipra\) proves
\[
 C(0,H(0;N,K)),\qquad
 n<N,\ C(n,D(n,k))\ \longrightarrow\ C(n+1,k)
\]
under a fixed context. Booleanize \(C\), hold \(N,K\) fixed, and
induct on
\[
             t\le N\ \longrightarrow\ C(t,H(t;N,K)).
\]
Lemma~\ref{lem:horizon} makes the previous observation exactly the
premise required by the local step. At \(t=N\), the terminal horizon
is \(K\), giving \(C(N,K)\). If the base theorem is available for
every demand and the retained context has no free \(N,K\), universal introduction over the independently fixed
\(N,K\) yields \(\forall N\forall K\,C(N,K)\).
The induction itself never uses \(\forall k\,C(t,k)\) as its invariant.

For example, if \(D(n,k)=2k+1\), elementary induction gives
\[
            H(t;N,K)=2^{N-t}(K+1)-1\qquad(t\le N).
\]
The resulting exponential precision demand is harmless for PR
definability. If \(D\) is nonmonotone, the original backward recurrence
still works; no maximum or monotonicity assertion is substituted for
its exact values. Several simultaneous queries are handled by the
bounded-prefix construction from the earlier packet, provided the
prefix validity interface is used explicitly.

\subsection{The new term-to-symbol compiler}

The repository represents a primitive-recursive function by an inductive
type \code{PR k}, with constructors for zero, successor, projections,
composition, and primitive recursion. A term \code{PTm k} may contain
variables and applications of these symbols. Define a compiler
\(S:\code{PTm k}\to\code{PR k}\) by
\[
 S(x_i)=\operatorname{proj}_i,
 \qquad
 S(f(t_0,\ldots,t_{m-1}))
   =\operatorname{comp}(f,S(t_0),\ldots,S(t_{m-1})).
\]

\begin{theorem}[Internal correctness of term compilation]
\label{thm:term-compile}
For every term \(t\), arbitrary tuple of substituted terms \(\sigma\),
and arbitrary context \(\Gamma\), there is an embedded arithmetic proof
\[
 \Gamma\vdash_{\mathrm{PRAAx}}
      \operatorname{app}(S(t),\sigma)=t[\sigma].
\]
The supplied Lean theorem \code{TermCompiler.symbol_correct} constructs
this proof and was checked by the Lean kernel.
\end{theorem}

\begin{proof}
Induct on the finite syntax of \(t\) in the metatheory. For a variable,
the projection axiom gives the equation. For an application, the
composition axiom gives
\[
 \operatorname{app}(S(t),\sigma)
   =f(\operatorname{app}(S(t_0),\sigma),\ldots,
       \operatorname{app}(S(t_{m-1}),\sigma)).
\]
The induction hypotheses are already object proofs for its argument
equations. Equality congruence and transitivity replace those arguments
by \(t_i[\sigma]\), yielding the desired equation. This proof consumes
only defining axioms and equality rules; no theorem about the truth of
arbitrary denotations is used.
\end{proof}

\subsection{The concrete compiler and the exact completed boundary}

In \code{DemandCompiler.lean}, the input types are
\code{D : PR 2}, \code{b : PR 1}, and \code{F : PR 3}; the output is
\code{compiled D b F : PR 2}. The intermediate symbols are:

\begin{center}
\begin{tabular}{lll}
\toprule
Symbol & Arity & Arguments\\
\midrule
\code{backwardSymbol D} & 3 & reverse count, terminal time, demand\\
\code{horizonSymbol D} & 3 & terminal time, demand, forward time\\
\code{forwardSymbol D b F} & 3 & forward count, terminal time, demand\\
\code{compiled D b F} & 2 & terminal time, demand\\
\bottomrule
\end{tabular}
\end{center}

The primitive-recursion convention is the actual upstream convention:
the recursion counter comes first, and a step symbol receives
\((\text{index},\text{accumulator},\text{parameters})\).
Predecessor and truncated subtraction are themselves explicitly built
from \code{PR.prec}. This avoids any unverified claim that a host
function is primitive recursive.

The checked theorem \code{denote_compiled} identifies the generated
syntax with the earlier semantic two-pass program. The checked host
theorems \code{compiled_zero} and \code{compiled_succ} establish the
desired recurrence equations for its denotation. In addition,
\code{object_backward_zero} and \code{object_backward_succ} construct
actual \code{Prf PRAAx} proofs of the generic backward pass equations
for arbitrary fixed \(D\). They use Theorem~\ref{thm:term-compile} to
handle the composite step body.

The full object-theory versions of
\eqref{eq:demand-zero}--\eqref{eq:demand-successor} are stated as
concrete target propositions in \code{ObjectTargets.lean}; they are
\emph{not yet implemented as Lean proofs}. Their human proof is
Theorem~\ref{thm:demand}. The remaining Lean work is finite and explicit:
derive the required subtraction equations internally, prove backward
coherence by equation induction, prove guarded horizon and forward
coherence, and unfold the final composition. Section~\ref{sec:lean}
records the exact order and the upstream APIs involved.

\subsection{The boundary with nested adaptive demands}

The theorem permits a changing parameter because \(D(n,k)\) does not
depend on the as-yet-uncomputed value \(X(n,\cdot)\). It does not
license an arbitrary nested recursion in which the next recursive
argument is obtained from a previous recursive answer. For a branching
request tree whose child parameters are fixed PR functions of the
parent and whose depth is \(N\), a supplied PR branching and parameter
envelope makes bounded bottom-up evaluation available. For a genuinely
adaptive algorithm, a fixed PR transition and the lexicographic
criterion in the next section provide another route. Without such a
resource argument, a total host definition or a familiar termination
proof does not establish PR implementability.
